\begin{claim}{D.72}{Proved}{The full mod-two cochain comparison for a pants involution}
Let $\alpha$ be the actual normalized representative of any
nonidentity finite-order element of the full pure three-cusp
loop-outer group. Retain all the data of D.34:
\[
 F=\widehat F(x,y),\qquad
 \alpha(x)=x^{-1},\quad \alpha(y)=f^{-1}y^{-1}f,
 \quad f\in F',\qquad E=F\rtimes_\alpha\langle s\rangle.
                                                               \tag{D.72a}
\]
Write $e_2\in\widehat\Z$ for the element whose $2$-adic
coordinate is one and whose other primary coordinates are zero.
The three actual section generators and their label maps are
\[
 s_0=s,\qquad s_1=x^{e_2}s,\qquad s_2=y^{e_2}fs,
 \qquad \iota_i:C_2\longrightarrow E,\quad\tau\longmapsto s_i.
                                                               \tag{D.72b}
\]
Let $\mathcal C^\bullet(E)$ be the differential graded
$\F_2$-algebra of normalized continuous group cochains with
trivial coefficients, with its ordinary cup multiplication.
Put
\[
 \mathcal A=\F_2[t_0,t_1,t_2]/(t_it_j:i\ne j),
              \qquad |t_i|=1,\qquad d_{\mathcal A}=0.
                                                               \tag{D.72c}
\]
The following is an actual surjective quasi-isomorphism of
unital differential graded algebras:
\[
 \begin{split}
 q:\mathcal C^\bullet(E)&\longrightarrow\mathcal A,\\
 q_0(a)&=a,\\
 q_n(c)&=\sum_{i=0}^2 c(s_i,\ldots,s_i)t_i^n
                                      \quad(n>0).
 \end{split}                                                   \tag{D.72d}
\]
Thus these are the complete mod-two cohomology ring and a
specified formality comparison, without an assumption that
the full centralizers of the sections are trivial.

Every sequence of at least three positive-degree cohomology classes whose
neighboring cup products are zero has a defined higher
Massey product containing zero. There is also the following
exact finite-matrix image calculation. For $n\geq1$, let
$U_{n+1}(\F_2)$ be the group of upper unitriangular matrices,
and define
\[
 \begin{gathered}
 \mathcal X_n=\bigl\{(\chi_1,\ldots,\chi_n)\in H^1(E,\F_2)^n:\
                    \chi_j\chi_{j+1}=0\ (1\leq j<n)\bigr\}.
 \end{gathered}
\]
The map reading the first superdiagonal gives a surjection
\[
 \operatorname{Hom}_{\rm cont}(E,U_{n+1}(\F_2))
                  \longrightarrow\mathcal X_n.
                                                               \tag{D.72e}
\]
It has exactly the displayed image. A representing
homomorphism can be chosen with the section values
\[
 s_i\longmapsto
       I+\sum_{j=1}^n\chi_j(s_i)E_{j,j+1}
                                  \quad(0\leq i<3).
                                                               \tag{D.72f}
\]
Writing $F_0=0,F_1=1,F_{k+2}=F_{k+1}+F_k$ for the
ordinary Fibonacci numbers, the number of possible
superdiagonal tuples is
\[
                              |\mathcal X_n|=F_{n+2}^{\,3}.
                                                               \tag{D.72g}
\]
No equality of full profinite groups, or solution of the
third-cusp conjugator equation, is inferred from these
cochain and matrix comparisons.
\end{claim}
\begin{proof}
\emph{The three literal restrictions and their normalizers.}
D.34 proves that the elements in (D.72b) have order two
and represent precisely the three $F$-conjugacy classes
of sections. The chosen generator of each quotient is
retained. Put $C_i=C_F(s_i)$ and $N_i=N_E(\langle s_i\rangle)$.
There are inverse continuous homomorphisms
\[
 C_i\times C_2\longrightarrow N_i,
       \quad(h,\tau^\epsilon)\longmapsto hs_i^\epsilon,
 \qquad
 w\longmapsto(ws_i^{-\epsilon},\tau^\epsilon),
                                                               \tag{D.72h}
\]
where $\epsilon$ is the actual image of $w$ in $E/F=C_2$.
A normalizer fixes the unique nonidentity element of its
order-two subgroup, so these formulas have the stated
domains and are inverse. Each $C_i$ is a closed subgroup
of the free profinite group $F$; consequently it has
$2$-cohomological dimension at most one. D.34 gives
$H^1(C_i,\F_2)=0$. Hence all its positive-degree
cohomology with these coefficients vanishes.

The Hochschild--Serre sequence for
$C_i\times C_2\to C_2$ has just its row zero. Inflation
and restriction to its displayed section therefore give
inverse isomorphisms
\[
           H^k(N_i,\F_2)\ \mathop{\rightleftarrows}^\sim\
                         H^k(C_2,\F_2)\qquad(k\geq0).
                                                               \tag{D.72i}
\]
The Brown restriction theorem
\cite[Theorem~2 and the finite-dimensional argument in
\S1.1]{A-KochloukovaZalesskii2007} applies to $E$:
$F$ is open of cohomological dimension one, and every
finite subgroup injects into $E/F$, so $E$ has no
$C_2\times C_2$. The two-row spectral sequence for
$F\to E\to C_2$ bounds $\dim H^k(E,\F_2)$ by three
for $k>1$. Its dense restriction embedding is therefore
onto its Hausdorff target. With the three representatives
and (D.72i), it gives the actual isomorphisms
\[
 (\iota_0^*,\iota_1^*,\iota_2^*):H^k(E,\F_2)
                   \xrightarrow{\ \sim\ }\F_2^3
                                      \qquad(k>1).
                                                               \tag{D.72j}
\]

\emph{Degree one and cochains in every degree.}
Let $u:E\to\F_2$ be the quotient character. Let
$\xi_x,\xi_y:E\to\F_2$ read the two mod-two
abelianization coordinates on $F$, and be zero on $s$.
These are actual characters because $\alpha$ is
minus identity on the full abelianization and therefore
identity modulo two. Every character of $E$ is
uniquely determined by its values on $x,y,s$, and
each triple of values is permitted by the semidirect
relations. Thus $u,\xi_x,\xi_y$ form a basis, with
the exact restriction table
\[
 \begin{array}{c|ccc}
       &s_0&s_1&s_2\\ \hline
 u     &1&1&1\\
 \xi_x &0&1&0\\
 \xi_y &0&0&1.
 \end{array}                                                   \tag{D.72k}
\]
Its inverse basis is
$\eta_0=u+\xi_x+\xi_y,\eta_1=\xi_x,\eta_2=\xi_y$;
each $\eta_i$ restricts to one on $s_i$ and zero
on the other two generators. Consequently (D.72j)
also holds in degree one. Degree zero is the common
scalar field, as in (D.72d).

For an order-two group with generator $\tau$, a
normalized cochain in degree $k>0$ is determined by
its value on $(\tau,\ldots,\tau)$. The differential
is zero: all internal terms have an identity argument,
and the two endpoint terms cancel in characteristic
two. The cup product multiplies the two specified
values. Its normalized cochain algebra is therefore
literally $\F_2[t]$, with $t$ in degree one and zero
differential. The three restrictions agree in degree
zero and give exactly (D.72d), including its cup
multiplication and its unit. This proves that $q$
is a differential graded algebra homomorphism.
The cohomology calculations above prove that it is
a quasi-isomorphism.

Surjectivity is also literal on cochains. Define the
finite quotient character
\[
 \rho=(u,\xi_x,\xi_y):E\longrightarrow\F_2^3,
 \qquad
 v_0=(1,0,0),\quad v_1=(1,1,0),\quad v_2=(1,0,1).
                                                               \tag{D.72l}
\]
For $k>0$, the continuous normalized cochain
\[
 e_{i,k}(g_1,\ldots,g_k)=
   \begin{cases}
       1,&\rho(g_1)=\cdots=\rho(g_k)=v_i,\\
       0,&\text{otherwise}
   \end{cases}
                                                               \tag{D.72m}
\]
satisfies $q(e_{i,k})=t_i^k$. The values $v_i$
are distinct and nonzero, so this formula is
normalized and gives all three basis vectors.
It factors through the stated finite quotient;
no extension of a function from an unspecified
subset is needed. Therefore
$\mathcal K^\bullet=\ker q$ is an acyclic differential
graded ideal, by the cohomology exact sequence of
this degreewise surjection.

The target is the graded algebra associated with
the eight-element Boolean ring $\F_2^3$ in
\cite[Definition~6.1]{R-PalQuick2025}. The displayed
restriction is a direct formality proof for these
particular extensions; it does not apply the
real-projective hypothesis to $E$.

\emph{All defining systems and their continuous lifts.}
Take $n\geq3$ positive-degree classes $a_1,\ldots,a_n$, with
$a_j\in H^{d_j}(E,\F_2)$, whose neighboring products
are zero. Their images in $\mathcal A$ have the
same property, by the proved cohomology isomorphism.
Choose cocycles $c_{j,j+1}$ representing them.
Their $q$-values are exactly the specified images,
because the target differential is zero.

We now choose cochains $c_{i,j}$ for every gap
$j-i\geq2$, including $(i,j)=(1,n+1)$, with
\[
 |c_{i,j}|=\sum_{k=i}^{j-1}d_k-(j-i-1),\qquad
 q(c_{i,j})=0,\qquad
 dc_{i,j}=\sum_{i<k<j}c_{i,k}\smile c_{k,j}.
                                                               \tag{D.72n}
\]
Suppose all smaller gaps have been chosen.
The right side is a cocycle: applying $d$ and
using the equations for smaller gaps makes every
triple product appear twice, so the terms cancel.
Its $q$-image is zero. At gap two this is exactly
the assumed neighboring product. At a larger gap,
each summand has a factor from a previously chosen
higher entry with $q$-value zero. The cocycle is
therefore in $\mathcal K^\bullet$. Acyclicity gives
the required $c_{i,j}$ inside that kernel. All its
degrees are positive because every $d_j\geq1$.
All cochains in this construction are continuous,
since the complex and its kernel consist of
continuous cochains throughout.

Omitting the final entry gives a defining system
for the higher Massey product. Including it in
(D.72n) exhibits its value as a coboundary. Thus
the product is defined and contains zero, with
every degree and every actual lifting complex
specified.

\emph{The exact upper unitriangular image.}
For a tuple in $\mathcal X_n$, apply the same
construction with $d_j=1$ and $c_{j,j+1}=\chi_j$.
The gap recursion also applies for $n=2$: its single
higher entry is obtained from acyclicity of
$\mathcal K^\bullet$ in exactly the displayed degree.
All entries are now normalized one-cochains.
Set
\[
   M(g)_{i,i}=1,\qquad M(g)_{i,j}=c_{i,j}(g) (i<j),
                   \qquad M(g)_{i,j}=0 (i>j).
                                                               \tag{D.72o}
\]
For a one-cochain the equation in (D.72n) reads
\[
 c_{i,j}(gh)=c_{i,j}(g)+c_{i,j}(h)
                     +\sum_{i<k<j}c_{i,k}(g)c_{k,j}(h).
\]
These are exactly all entries of $M(gh)=M(g)M(h)$.
The finite matrix-valued map is continuous and
has the required superdiagonal. Its higher
entries vanish on every $s_i$ because they lie
in $\ker q$, proving (D.72f). Conversely any
continuous unitriangular representation has
\(d c_{j,j+2}=\chi_j\smile\chi_{j+1}\), so its
superdiagonal belongs to $\mathcal X_n$.
This proves both directions of the image statement.
For $n=1$ the map is just the character itself.

Finally use the isomorphism (D.72k) to record
each $\chi_j$ by its three section values.
The neighboring cup product vanishes exactly
when those two three-bit vectors have disjoint
supports, by (D.72j). Each of the three coordinates
therefore gives a binary string of length $n$
with no two consecutive ones. Their number
$b_n$ satisfies $b_0=1,b_1=2$ and
$b_n=b_{n-1}+b_{n-2}$ by its final digit.
Hence $b_n=F_{n+2}$, and the three independent
coordinates give (D.72g).

The same three restriction maps for the standard
involution $\alpha_0(x)=x^{-1},\alpha_0(y)=y^{-1}$
give the explicit zigzag
\[
 \mathcal C^\bullet(E)\xrightarrow{\ q\ }\mathcal A
       \xleftarrow{\ q_{\rm std}\ }
       \mathcal C^\bullet(F\rtimes_{\alpha_0}C_2).
                                                               \tag{D.72p}
\]
Both arrows are the specified quasi-isomorphisms
of differential graded algebras. This is the
exact comparison established here. It leaves
the third-cusp equation of D.35 and the full
section-centralizer equations unchanged.
\end{proof}
